% Partie : applied-modal-logics
% Chapitre : epistemic-logic
% Section : public-announcement-logic-semantics

\documentclass[../../../include/open-logic-section]{subfiles}

\begin{document}

\olfileid[fr]{aml}{el}{psm}

\olsection{Sémantique de la logique des annonces publiques}

Les modèles relationnels sont les mêmes qu'en logique épistémique.
En revanche, la sémantique de l'opérateur d'annonce publique est nouvelle.

\begin{defn}\ollabel{defn:mmodels}
  La \emph{vérité d'!!a{formula}~$!A$ en~$w$} dans
  un modèle~$\mModel M = \tuple{W, R, V}$, notée
  $\mSat{M}{!A}[w]$, est définie inductivement comme suit :
  \begin{enumerate}
  \tagitem{prvFalse}{\indcase{!A}{\lfalse}{On n'a jamais $\mSat{M}{\lfalse}[w]$}.}{}
  \tagitem{prvTrue}{\indcase{!A}{\ltrue}{On a toujours $\mSat{M}{\ltrue}[w]$}.}{}
  \item $\mSat{M}{p}[w]$ si et seulement si $w \in V(p)$.
  \tagitem{prvNot}{\indcase{!A}{\lnot !B}{$\mSat{M}{\indfrm}[w]$ si et seulement si
    $\mSat/{M}{!B}[w]$}.}{}
  \tagitem{prvAnd}{\indcase{!A}{(!B \land !C)}{$\mSat{M}{\indfrm}[w]$ si et seulement si
    $\mSat{M}{!B}[w]$ et $\mSat{M}{!C}[w]$}.}{}
  \tagitem{prvOr}{\indcase{!A}{(!B \lor !C)}{$\mSat{M}{\indfrm}[w]$ si et seulement si
    $\mSat{M}{!B}[w]$ ou $\mSat{M}{!C}[w]$} (éventuellement les deux).}{}
  \tagitem{prvIf}{\indcase{!A}{(!B \lif !C)}{$\mSat{M}{\indfrm}[w]$ si et seulement si
    $\mSat/{M}{!B}[w]$ ou $\mSat{M}{!C}[w]$}.}{}
  \tagitem{prvIff}{\indcase{!A}{(!B \liff !C)}{$\mSat{M}{\indfrm}[w]$ si et seulement si
    soit on a à la fois $\mSat{M}{!B}[w]$ et $\mSat{M}{!C}[w]$, soit
    on n'a ni $\mSat{M}{!B}[w]$ ni $\mSat{M}{!C}[w]$}.}{}
  \item\ollabel{defn:sub:mmodels-box}
    \indcase{!A}{\Knows_a !B}{$\mSat{M}{\indfrm}[w]$ si et seulement si
    $\mSat{M}{!B}[w']$ pour tout $w' \in W$ tel que $R_a ww'$}
  \item\ollabel{defn:sub:mmodels-pal}
    \indcase{!A}{[!B] !C}{$\mSat{M}{\indfrm}[w]$ si et seulement si
    $\mSat{M}{!B}[w]$ entraîne $\mSat{M \mid !B}{!C}[w]$}

  où la restriction $\mModel M \mid !B = \tuple{W', R', V'}$ est
  définie comme suit :
  \begin{enumerate}
    \item $W' = \Setabs{ u \in W}{\mSat{M}{!B}[u]}$. Les mondes de
      $\mModel M \mid !B$ sont donc les mondes de~$\mModel M$ où $!B$
      est vraie.
    \item $R'_a = R_a \cap (W' \times W')$. La relation d'accessibilité
      de chaque agent est restreinte aux mondes conservés dans~$W'$.
    \item $V'(p) = \Setabs{ u \in W'}{u \in V(p) }$. Les valuations
      propositionnelles restent inchangées aux mondes conservés :
      cette mise à jour informationnelle ne modifie pas la valeur de
      vérité des variables propositionnelles en ces mondes.
  \end{enumerate}
  \end{enumerate}
\end{defn}

Si aucun monde ne satisfait l'annonce, la restriction est vide.
La clause d'annonce n'y évalue toutefois aucune formule : son
antécédent est alors faux au monde de départ. Si l'on ajoute
l'opérateur de connaissance commune, on conserve sa clause
sémantique précédente, appliquée au modèle courant, y compris
aux modèles restreints.

La particularité de la logique des annonces publiques est donc que
la vérité d'!!a{formula} en un monde de~$\mModel M$ peut parfois
nécessiter une évaluation dans un autre modèle que~$\mModel M$.

La sémantique traite l'opérateur d'annonce comme un opérateur de
type~$\Box$. Ainsi, si !!a{formula}~$!A$ ne peut être annoncée
de manière véridique en un monde, $[!A]B$ y est vraie par vacuité,
comme toutes les !!{formula}s en $\Box$ le sont dans un monde
sans successeur.

\begin{figure}
  \begin{center}
    \begin{tikzpicture}[modal]
      \node[world] (w1) [label={below left:$p, \lnot q$}]{$w_1$};
      \node[world] (w2) [left=of w1, label={left:$\lnot p, \lnot q$}]{$w_2$};
      \node[world] (w3) [above=of w1, label={right:$p, q$}] {$w_3$};
      \draw[<->] (w1) to [swap] node {$b$} (w2);
      \draw[->,loop below] (w1) to node {$a, b$} (w1);
      \draw[<->] (w1) to [swap] node {$a$} (w3);
      \draw[->,loop above] (w2) to node {$a, b$} (w2) ;
      \draw[->,loop above] (w3) to node {$a, b$} (w3) ;

      \node[world](v1)[right of=w1, xshift=1.25in, label={right:$p, \lnot q$}]{$w'_1$};
      \node[world](v3)[above of=v1,label={right:$p,q$}]{$w'_3$};
      \draw[<->] (v1) to node {$a$} (v3);
      \draw[->,loop below] (v1) to node {$a,b$} (v1);
      \draw[->,loop above] (v3) to node {$a,b$} (v3) ;

      \node(m)[below=of w1]{$\mModel M$};
      \node(m')[below=of v1]{$\mModel M \mid p$};

      \draw[-,dotted] (w1) to node {\footnotesize annonce de $p$}(v1) ;
    \end{tikzpicture}
  \end{center}
  \caption{Avant et après l'annonce publique de $p$.}
  \ollabel{fig:announcement-example}
\end{figure}

On peut voir une annonce publique d'!!a{formula} comme une
réduction du modèle aux mondes où la !!{formula} était vraie.
La \olref{fig:announcement-example} montre l'effet de l'annonce
publique de~$p$. Les primes y désignent simplement les mêmes mondes
après restriction. Dans cet exemple, l'agent~$b$ apprend que~$p$
à la suite de l'annonce, tandis que l'agent~$a$ ne l'apprend pas à
cette occasion, puisque $a$ savait déjà que~$p$ était vraie.

Plus formellement, on a $\mSat{M}{\lnot \Knows_b p}[w_1]$, mais
$\mSat{M
\mid p}{\Knows_b p}[w'_1]$. On en déduit
$\mSat{M}{[p] \Knows_b
p}[w_1]$. On peut même dire davantage sur le résultat de l'annonce :
on a $\mSat{M}{[p]\CKnows_{\{a,b\}} p}[w_1]$. Autrement dit,
après l'annonce, $p$ devient une \emph{connaissance commune}.

Mais est-ce toujours le cas ? Une annonce véridique de~$!A$
fait-elle \emph{toujours} de $!A$ une connaissance commune ? La
réponse, peut-être surprenante, est non. Il est même possible
d'annoncer de manière véridique des !!{formula}s qui cessent
d'être vraies après leur annonce. Considérons, par exemple,
l'annonce de $p \land \lnot \Knows_b p$ en~$w_1$ dans la
\olref{fig:announcement-example}. Les modèles $\mModel M \mid p$ et
$\mModel M \mid (p \land \lnot \Knows_b p)$ sont ici différents :
le premier conserve aussi le troisième monde, tandis que le second
ne conserve que le premier.\footnote{La source affirme à tort que ces deux modèles restreints sont égaux. Le troisième monde satisfait l'annonce simple, mais pas l'annonce plus forte. Cette inégalité ne change pas le contre-exemple qui suit.} Dans chacun, l'agent b sait que p au
premier monde conservé. Nous avons déjà établi que
$\mSat{M \mid p}{\Knows_b
p}[w'_1]$. L'annonce plus forte ne laisse, elle aussi, qu'un monde
accessible à b où p est vraie. On a donc
$\mSat{M \mid (p \land \lnot \Knows_b p)}{\lnot
(p \land \lnot \Knows_b p)}[w'_1]$ : cette !!a{formula} devient
fausse après son annonce.

\end{document}
